\section{Introduction}\label{sec:intro}

Since the release of OpenAI's ChatGPT\footnote{https://chatgpt.com/} in November 2022, chatbots based on Large language models (LLMs) have become ubiquitous in the public consciousness.
Due to their continued improvement, going from simple chat interfaces to fully integrated agentic systems like Anthropic's Claude Code\footnote{https://claude.com/product/claude-code}, they are increasingly helpful in any workplace dealing with writing, programming, or project planning.
One such workplace is academia, and it is not surprising that researchers across disciplines have begun to adopt LLMs for editing or writing papers, peer reviews, or grant applications \cite{Kwon2025}.
Proponents of the technology praise it as a tool for accessibility and inclusivity \cite{Bao2025_linguistic, Berdejo-Espinola2023, Prakash2025}, arguing that it enables non-native English speakers (NNES) to overcome language barriers by translating or performing grammar checks \cite{Kousha2025}.
However, the surge of ChatGPT's popularity is accompanied by concerns of misuse \cite{Liebrenz2023, Zheng2023}.
LLMs are prone to hallucination, so if their output is not thoroughly checked before using it in a paper, false or misleading statements might be published, especially if the review is also delegated to an LLM \cite{Alansari2026, Lindsay2023, Walters2023, Xu2025}.
Chatbots also make scientific fraud much easier, with the possibility of fabricating entire papers at the cost of a single prompt \cite{Kendall2024}.  
Reviewers relying on LLMs are vulnerable to hidden prompts in the paper's text, telling the LLM to give a positive review \cite{Lin2025_hidden_prompts}.
In light of these concerns, many journals have issued guidelines about the disclosure of LLM-use \cite{Brainard2023}. 
However, the perception of LLM-use leading to inaccuracies in the text, together with journals restricting or explicitly banning certain use-cases may incentivize academics against disclosing \cite{Thorp2023}.
Consequently, when trying to investigate the prevalence of LLM-use in academia, and especially in the paper-writing process, relying on self-disclosure statements will lead to underestimation \cite{Kousha2024, Kwon2025}.
\\
Regardless of the ethical concerns about using LLMs as paper-writing assistants, their prolonged use in academic text production will have an impact on academic writing as a whole.
LLMs have a distinct writing style.
Sentences generated by them have been shown to be less diverse in length \cite{Desaire2023}, but to have higher lexical diversity and positive sentiment \cite{Mak2025}.
Famously, they also show a preference for certain words, an observation which has lead to much ridicule around the most popular example ``to delve'' \cite{MacClure2024}.
The reason for this behavior is not entirely clear.
Juzek and Ward \cite{Juzek2025} replicated the training process of chatbots and tried to artificially create word preferences in the model.
They found no evidence to suggest that the training data, model architecture, or training algorithm have an influence in forming these idiosyncratic word choices.
However, the authors were able to induce an LLM with preferred words during reinforcement learning with human feedback (RLHF), the continued training scheme behind the success of modern Chatbots \cite{Bai2022}.
During RLHF, humans provide learning goals by rating the interaction with the model.
For commercial Chatbots, this work is often performed by precariously employed workers in the global south, for example in Nigeria, where ``delve'' is used much more frequently than in other English-speaking countries \cite{Perrigo2023, Hern2024}.
\\
These selective word preferences can be leveraged to estimate the proportion of LLM-processed text in a corpus.
Kobak et al. \cite{Kobak2025} track changes in word frequencies in biomedical abstracts from the PubMed database.
They compare a word's observed frequency in a given year with a projection based on its frequencies two and three years back.
The metric obtained by subtracting the predicted from the observed frequency is called excess frequency: A change in frequency that diverges from previous trends.
They identify many words that appear with excess frequency in abstracts from 2024.
This is notable, because for previous years, the observed frequency gaps are much smaller.
Because of that, a threshold can be applied to excess frequency, such that most words from pre-2024 don't exceed it, but many from 2024 do.
In addition, the majority of pre-2024 excess words are clearly related to a specific research topic, such as ``ebola'' in 2015, ``convolutional'' in 2017, or ``coronavirus'' in 2020.
In contrast, many excess words in 2024 are content-independent style words -- ``underscore'', ``potential'', ``significant'' or ``pivotal'', to name just a few.
The qualitative and quantitative differences between excess words from before and after 2024 point to a drastic shift in writing patterns that can only be explained by the adoption of LLMs into the writing process.
The identified excess style words for 2024 can therefore be interpreted as marker words for LLM-use.
It's important to note that the marker words should be applied to analyses on the corpus level and not to individual documents -- one ``delve'' does not an LLM-generated text make. 
While these words appear more frequently in LLM outputs, some of them still have very low frequency in both human and machine text.
There are texts by human authors that include them, and LLM-generated texts that don't.
\\
Using the combined excess frequencies of a selection of excess style words, Kobak et al. derive a lower bound for LLM-use estimates.
Their reasoning is quite intuitive: all instances of using one of the relevant words that go beyond the predicted frequency must have been caused by LLM-use.
Excess frequency as a metric has the advantage of not needing to distinguish between using an LLM write a whole paragraph, or suggesting changes on previously written text.
Both use-cases will lead to some LLM-suggested words being present in the final product.
\\
Here, we expand on their work on excess word frequency as a lower bound estimate of LLM-use in biomedical abstracts by making the following contributions:
\begin{itemize}
    \item \textbf{Widening the focus} from an analyzing only the abstract to the entire paper. By using the open-access PubMed Central (PMC \cite{pmc}) dataset  instead of PubMed, the proportion of LLM-assisted text in the main text of the paper can be compared with usage in the abstract, and different sections within the body can be compared against each other. 
    \item Deriving a \textbf{stronger lower bound} metric going beyond excess frequency. We show that while excess frequency is an intuitive estimate for LLM-use, it severely underestimates the true proportion of LLM-assisted text. 
    \item Computing \textbf{uncertainty measures} and performing tests on \textbf{simulated data} to increase confidence in the estimates. A disadvantage of the previous approach is the lack of a ground-truth set to test performance with. By simulating the data generation process in a Binomial model of word occurrence, we can evaluate our metric for several ground-truth levels of LLM-use.
    \item \textbf{Comparing} the LLM-use estimates with those of other estimation approaches, as well as survey data asking researchers about their LLM-use. Unlike previous studies, we find that LLM-like language has almost completely permeated academic writing.
\end{itemize}

This thesis is structured in the following way: In Section \ref{sec:related}, other approaches to estimate LLM-use in academia are discussed, such as comparable word frequency measures, word distribution modeling and LLM-detectors.
Section \ref{sec:methods} gives an overview of the data processing pipeline, the metrics used to estimate LLM-use, as well as the validation procedure employed to increase confidence in the estimates attained.
This is followed by Section \ref{sec:results}, where our estimates are reported and compared to results from other work.
Finally, the results, along with their impact and limitations are discussed in Section \ref{sec:discussion}.
\\
All code is available on GitHub\footnote{https://github.com/LenaHolzwarth/llms\_in\_academia}.
%The PMC dataset can be downloaded here\footnote{https://ftp.ncbi.nlm.nih.gov/pub/pmc/deprecated/oa\_bulk/}.